\subsection{Class two for every fixed finite puncture set}

\begin{statement}{A.50}{Proved}{Every ordinary finite-puncture ghost lifts to full profinite class two}
Let $S=\{c_1,\ldots,c_r\}\subset\Q$ have distinct elements, with
$r\ge2$, and put
\[
 V_S=\mathbb A^1_\Q\setminus S,\qquad
 E_S=\pi_1(V_S),\qquad \Pi_S=\pi_1(V_{S,\bar\Q}).
\]
Let $U_i\in\widehat{\Q^\times}$ and, for every finite prime $p$,
let $x_p\in\Q_p\setminus S$ satisfy
\[
                (U_i)_p=\kappa_p(x_p-c_i)\quad(1\le i\le r).
\tag{A.50.1}
\]
There is a continuous section of
\[
 E_S/\Gamma_3\Pi_S\longrightarrow G_\Q
\tag{A.50.2}
\]
whose pullbacks of the unit classes $\kappa(t-c_i)$ are the
specified $U_i$.  This is the actual full profinite class-two
quotient, with all coefficient primes retained.  No local
section theorem is used.

For the classes of A.48, none of these sections is the image
of a rational point section or of a cuspidal section.  Thus
ordinary infinite-support abelian data survive to this quotient
for every fixed finite rational puncture set.  The statement
concerns the global lifting problem; it does not prescribe
the central coordinates of the local point sections.
\end{statement}
\begin{proof}
Choose $b\in\Q\setminus S$ as a rational base point.  The unit
maps $t-c_i:V_S\to\mathbb G_m$ give the canonical coordinate
isomorphism
\[
 T_S:=\Pi_S^{\rm ab}\xrightarrow{\ \sim\ }
                      \widehat\Z(1)^r.
\tag{A.50.3}
\]
One can check the isomorphism on geometric meridians: a meridian
at $c_i$ has Kummer degree one for $t-c_i$ and zero for the
other units, and the meridians at the finite punctures freely
generate the geometric profinite group.  Rationality of the
punctures gives the cyclotomic action on every coordinate.
The point $b$ splits the arithmetic extension.  Hence the
specified $U_i$ define an abelian section $s_1$ of
$E_S^{\rm ab}:=E_S/\Gamma_2\Pi_S$: its coordinate cocycles,
relative to this splitting, represent
$U_i-\kappa(b-c_i)$ in additive Kummer notation.

The kernel of the next projection is the explicit module
\[
 M_S:=\Gamma_2\Pi_S/\Gamma_3\Pi_S
    \xrightarrow{\ \sim\ }\bigwedge^2 T_S
    =\bigoplus_{i<j}\widehat\Z(2)e_{ij}.
\tag{A.50.4}
\]
Here $e_{ij}$ is the commutator of the meridians at $c_i,c_j$,
in that order.  We recall a direct verification which also
supplies a continuous set-theoretic section of the projection.
The free class-two profinite group on $r$ generators has
coordinates
\[
 (a,h)\in\widehat\Z^r\times\widehat\Z^{\binom r2},\qquad
 (a,h)(a',h')=
  \bigl(a+a',\ h+h'+(a_i a'_j)_{i<j}\bigr).
\tag{A.50.5}
\]
The displayed multiplication is associative because its central
term is bilinear.  Its standard generators have independent
commutators $e_{ij}$.  In any profinite group of class at most
two, collecting words gives the reverse-ordered generator
powers followed by the commutator powers in (A.50.5).
The same identities extend from integral exponents to
profinite exponents by continuity.  The resulting map from
(A.50.5) to any such group is the unique continuous homomorphism
with the specified generator images.  Conversely the standard
coordinate group itself separates all these powers.
This proves the universal description and (A.50.4).
In particular $(a,0)$ is a continuous lift of each abelian
coordinate.  The arithmetic action on $M_S$ is $\chi^2$,
because it is the exterior square of (A.50.3); inner geometric
conjugation acts trivially on this central geometric subgroup.

Pulling back $E_S/\Gamma_3\Pi_S\to E_S^{\rm ab}$ along $s_1$
therefore defines an obstruction
\[
            o(s_1)\in H^2_{\rm cts}(G_\Q,M_S).
\tag{A.50.6}
\]
More explicitly, use $(a,0)$ and the point-$b$ splitting to
lift $s_1$ continuously; the defect of multiplication is a
continuous $M_S$-valued two-cocycle.  A continuous primitive
changes the lift into a section, and a section supplies such
a primitive.  This proves the obstruction assertion for
continuous cochains, without an inverse-limit substitution.

For each pair $i<j$, the open immersion which forgets the
other punctures, followed by an affine change of coordinate,
is the explicit $\Q$-morphism
\[
 f_{ij}:V_S\longrightarrow U,\qquad
 z=\frac{t-c_i}{c_j-c_i},\qquad
 1-z=\frac{-(t-c_j)}{c_j-c_i}.
\tag{A.50.7}
\]
It gives continuous homomorphisms of the arithmetic fundamental
groups and of their abelian and class-two quotients.  Use the
image base point $f_{ij}(b)$ and compatible geometric paths
to compare with the generators in A.38.  On (A.50.4), the
induced map is precisely
\[
 M_S\longrightarrow\widehat\Z(2),\qquad
          e_{ij}\longmapsto1,\quad e_{kl}\longmapsto0
                          \quad\text{for }(k,l)\ne(i,j).
\tag{A.50.8}
\]
Indeed the only surviving meridians are those of $c_i,c_j$;
their respective Kummer coordinates in (A.50.7) are $(1,0)$
and $(0,1)$.  Changes of geometric paths induce inner
conjugations, which act trivially on this commutator module.
Thus (A.50.8) also fixes its sign and its Galois action.

The image of $s_1$ in this two-puncture abelian quotient has
unit classes
\[
 A_{ij}=U_i/\kappa(c_j-c_i),\qquad
 B_{ij}=\kappa(-1)U_j/\kappa(c_j-c_i),
\tag{A.50.9}
\]
where the products in (A.50.9) are multiplicative.  At every
finite $p$, they are the ordinary unit-equation pair
\[
 \left(\frac{x_p-c_i}{c_j-c_i},\,
       \frac{c_j-x_p}{c_j-c_i}\right).
\]
Neither coordinate is zero and their sum is one.  A.27 and
A.38 give a continuous section of the two-puncture class-two
quotient with these abelian classes.  Naturality of the defect
cocycle under (A.50.7) shows that its obstruction is the
$(i,j)$ coordinate of (A.50.6).  It therefore vanishes.
Continuous cochains commute with this finite direct product
of coefficient modules, so the coordinate maps give an
isomorphism
\[
 H^2_{\rm cts}(G_\Q,M_S)
       \simeq\prod_{i<j}H^2_{\rm cts}
                        (G_\Q,\widehat\Z(2)).
\]
Every component of (A.50.6) is zero.  Its continuous primitive
constructs (A.50.2), as required.

Finally, in A.48 every $U_i$ has infinite support, and every
coordinate has infinitely many new support primes absent from
the other coordinates.  For a rational point $t=b'$, the
class of every $t-c_i$ is a rational number and has finite
support.  For a section in the packet of $c_i$, every unit
$t-c_j$ with $j\ne i$ extends across that cusp with the
nonzero value $c_i-c_j$.  Its pullback is therefore the
rational class $\kappa(c_i-c_j)$, which again has finite
support.  At infinity use the local parameter $z=1/t$:
\[
                      t-c_i=z^{-1}(1-c_i z).
\]
The principal-unit factor has compatible roots with value
one, so all the Kummer classes there have the same tangent
parameter with exponent $-1$.  They are equal to one another.
The distinct new-prime supports in A.48 exclude this last
possibility.  These comparisons all factor through
(A.50.3), so they exclude the indicated images in the actual
class-two quotient.
\end{proof}

\subsection{What one fixed finite quotient can test}

\begin{statement}{A.51}{Proved}{Infinite-support ghosts survive every fixed finite geometric quotient with full abelianization}
Let $N$ be a closed normal subgroup of $E_U$ contained in
$\Pi$, and suppose $Q=\Pi/N$ is finite.  Write
\[
 E_Q=E_U/N,\quad E_{Q^{\rm ab}}=E_Q/[Q,Q],\quad
 E_1=E_U/\Gamma_2\Pi.
\]
The natural map is an isomorphism of profinite extensions of
$G_\Q$,
\[
 E_U/(N\cap\Gamma_2\Pi)
       \xrightarrow{\ \sim\ }
                 E_Q\mathop{\times}_{E_{Q^{\rm ab}}}E_1.
\tag{A.51.1}
\]
For every such $Q$, the extension on the left has a continuous
section whose full abelian Kummer coordinates form an
ordinary infinite-support ghost of the form A.44.  Its image
in $E_Q$ is the image of the rational point section at $3$.
For perfect $Q$, every ordinary ghost admits a section of
this extension with that prescribed $E_Q$ image.

These sections are excluded from the rational-point and
cuspidal images by their full abelian coordinates.  Thus no
single finite geometric quotient, adjoined to the full
abelianization in the exact sense (A.51.1), eliminates all
ordinary infinite-support ghosts.  This includes any fixed
finite modular quotient with geometric group
$A_4,(\Z/2)^2,A_5$ or $\mathrm{PSL}_2(\mathbf F_7)$.
The assertion prescribes a global quotient section; matching
its finite-quotient localizations to the points $x_p$ is an
additional requirement.
\end{statement}
\begin{proof}
Put $T=\Pi^{\rm ab}$.  The map
\[
               \Pi\longrightarrow Q\times_{Q^{\rm ab}}T
\tag{A.51.2}
\]
is onto.  To verify this directly, start with $(q,t)$ in the
fibre product and lift $q$ to $g\in\Pi$.  The difference
between $t$ and the abelian image of $g$ lies in the kernel
of $T\to Q^{\rm ab}$.  This kernel is the image of $N$.
Indeed an element of $\Pi$ whose $Q$ image belongs to
$[Q,Q]$ can be multiplied by an element of $[\Pi,\Pi]$
with the inverse image in $[Q,Q]$, leaving an element of $N$
with the same abelian image.  Here the map of commutator
groups is onto, since every commutator in $Q$ lifts and $Q$
is finite.  Multiply $g$ by the corresponding element of
$N$ to obtain the desired pair.  The kernel of (A.51.2) is
$N\cap\Gamma_2\Pi$.  The same argument on each fibre over
$G_\Q$ proves surjectivity in (A.51.1); its kernel is the
same subgroup.  The resulting continuous bijection between
compact Hausdorff groups is a topological isomorphism.

Let $d$ be the exponent of $Q^{\rm ab}$, with $d=1$ when
$Q^{\rm ab}$ is trivial.  The recursive proof of A.44 works
with
\[
                 n_j=(j+k)!,\qquad k\ge4,
\tag{A.51.3}
\]
in place of $(j+4)!$.  We spell out the only indexed checks.
Every $n_j$ is divisible by $24$ and by $8$; a prime divisor
of $n_{j+1}$ is at most $j+k+1\le(n_j+1)^4$; and the
unchanged precision
$H_j(p)=2v_p(n_{j+1})+j+3$ is strictly larger than
$2v_p(n_j)$ and gives the required displacement at least
$v_p(n_{j+1})$.  All field, counting and new-prime arguments
in A.44 are unchanged algebraic identities with these $n_j$.
The construction therefore terminates at every stage and
has the same infinite, disjoint new-prime supports.

Choose $k$ so that $d\mid k!$.  The resulting classes
satisfy
\[
           A\equiv\kappa(3),\qquad
           B\equiv\kappa(-2)
                 \pmod{(\widehat{\Q^\times})^d},
\tag{A.51.4}
\]
meaning equality in $\Q^\times/(\Q^\times)^d$ under the
canonical finite-power projection.  Equivalently, relative
to the point-$3$ splitting of $E_1$, their abelian cocycle
has zero cohomology class after reduction in $T/dT$.
A coboundary representative in $T/dT$ lifts to an element
of $T$, so a geometric conjugation makes that reduced
cocycle literally zero.  Since $T\to Q^{\rm ab}$ factors
through $T/dT$, the resulting abelian section and the
point-$3$ section have exactly the same image in
$E_{Q^{\rm ab}}$.

Pair this representative of the abelian ghost section with
the point-$3$ section of $E_Q$.  Their equal images define
a continuous section of the fibre product in (A.51.1).
The isomorphism there gives the claimed section of the
actual quotient of $E_U$.  When $Q$ is perfect,
$E_{Q^{\rm ab}}=G_\Q$ and there is no matching condition;
the same pairing works for every abelian ghost section.
The infinite support and the two permanent valuation
patterns in A.44 exclude point and cusp images exactly
as at the end of its proof.  Finally the argument used no
restriction on the finite group beyond being a geometric
quotient stable under the arithmetic action, so it applies
to each of the specified finite modular images with its
actual arithmetic extension.  This proves the precise
finite-level assertion without an assertion about local
matching, an infinite congruence quotient, or a compatible
family as the quotient varies.
\end{proof}
